% =====================================================================
%  CHƯƠNG 4. THIẾT KẾ VÀ XÂY DỰNG PHẦN MỀM MÔ PHỎNG
% =====================================================================
\chapter{THIẾT KẾ VÀ XÂY DỰNG PHẦN MỀM MÔ PHỎNG}
\label{chap:4}

% ---------------------------------------------------------------------
\section{Quy trình mô phỏng}
\label{sec:4.1}

Hình~\ref{fig:quytrinh} mô tả quy trình mô phỏng từ chuỗi ký tự đầu vào đến
chuỗi ký tự khôi phục.

\begin{figure}[H]
\centering
\begin{tikzpicture}[
  node distance=5mm and 7mm,
  khoi/.style={draw,rounded corners=2pt,align=center,minimum height=11mm,
               text width=3.3cm,font=\small,fill=gray!6},
  kenh/.style={khoi,fill=gray!22},
  mui/.style={-{Stealth[length=2.2mm]},thick}]
  \node[khoi] (a) {Chuỗi ký tự\\ đầu vào};
  \node[khoi,right=of a] (b) {Đổi sang byte,\\ tách khối 4 bit $u$};
  \node[khoi,right=of b] (c) {Mã hóa\\ $c = uG$ (7 bit)};
  \node[kenh,below=of c] (d) {Kênh BSC\\ $r = c \xor e$};
  \node[khoi,left=of d] (e) {Hội chứng $s = r\HT$,\\ sửa một bit};
  \node[khoi,left=of e] (f) {Lấy 4 bit đầu,\\ ghép thành ký tự};
  \draw[mui] (a) -- (b);
  \draw[mui] (b) -- (c);
  \draw[mui] (c) -- (d);
  \draw[mui] (d) -- (e);
  \draw[mui] (e) -- (f);
\end{tikzpicture}
\caption{Quy trình mô phỏng mã hóa, gây nhiễu, phát hiện và sửa lỗi}
\label{fig:quytrinh}
\end{figure}

\canviet{Mô tả từng bước của quy trình (4.1.1): ký tự $\to$ byte $\to$ hai khối
4 bit; mã hóa từng khối; truyền qua kênh; tính hội chứng, sửa lỗi; lấy 4 bit
đầu; ghép các khối thành byte và ký tự. Thay hình trên bằng sơ đồ chính thức
nếu có.}

% ---------------------------------------------------------------------
\section{Cấu trúc các module}
\label{sec:4.2}

\begin{table}[H]
\centering
\caption{Các module của chương trình}
\label{tab:module}
\begin{tabular}{|p{3.9cm}|p{11.6cm}|}
\hline
\textbf{Module} & \textbf{Chức năng}\\ \hline
\code{hamming.py}    & Mã hóa và giải mã một khối: \code{hamming\_encode},
                       \code{calculate\_syndrome}, \code{hamming\_decode},
                       kiểm tra đầu vào \code{validate\_bits}.\\ \hline
\code{text\_codec.py} & Đổi chuỗi ký tự $\leftrightarrow$ dãy bit, tách và ghép
                       khối 4 bit, gọi \code{hamming.py} cho từng khối.\\ \hline
\code{channel.py}    & Giả lập kênh truyền có nhiễu (mục~\ref{sec:4.4}).\\ \hline
\code{test\_hamming.py} & Bộ ca kiểm thử tự động (mục~\ref{sec:5.1}).\\ \hline
\end{tabular}
\end{table}

Trong mã nguồn, biến \code{data} ứng với vectơ $u$, \code{codeword} ứng với
từ mã $c$, \code{received} ứng với $r$ và \code{syndrome} ứng với $s$; chỉ số
Python đếm từ $0$ nên bit thứ $i$ của báo cáo là phần tử $i-1$. Mã nguồn đầy
đủ ở Phụ lục~A.

\canviet{Mô tả chi tiết từng hàm (đầu vào, đầu ra, ngoại lệ). Xác nhận tên
hàm của \code{text\_codec.py} và \code{channel.py} theo mã nguồn cuối cùng.}

% ---------------------------------------------------------------------
\section{Thuật toán lõi}
\label{sec:4.3}

Mã hóa một khối $u$ theo \eqref{eq:bitkiemtra}. Với vectơ nhận được
$r = (r_1\,\dots\,r_7)$, mỗi hàng của $H$ cho một thành phần của hội chứng
$s = r\HT$:
\begin{equation}\label{eq:s-thanhphan}
  s_1 = r_1 \xor r_2 \xor r_4 \xor r_5, \quad
  s_2 = r_1 \xor r_3 \xor r_4 \xor r_6, \quad
  s_3 = r_2 \xor r_3 \xor r_4 \xor r_7 .
\end{equation}
Bảy hội chứng khác $0$ chính là bảy cột của $H$; chương trình tra
Bảng~\ref{tab:hoichung} để tìm vị trí cần lật.

\canviet{Mã giả hoặc lưu đồ của \code{hamming\_encode} và \code{hamming\_decode};
ví dụ chạy tay (1011 $\to$ 1011010; lỗi bit 3 $\to$ 1001010; $s = (0,1,1) = h_3$;
sửa lại 1011010; giải mã 1011). Ghi rõ: khi có đúng hai lỗi, hội chứng luôn
khác 0 và trùng hội chứng của một lỗi đơn, nên chương trình luôn sửa nhầm; khi
từ ba lỗi, hội chứng có thể bằng 0.}

% ---------------------------------------------------------------------
\section{Mô hình nhiễu}
\label{sec:4.4}

Kênh truyền được mô phỏng bằng kênh đối xứng nhị phân (BSC, \emph{binary
symmetric channel}) với xác suất đảo bit $p$, $0 \le p \le 1$: mỗi bit của
từ mã bị đảo độc lập với xác suất $p$. Vectơ nhận được là $r = c \xor e$,
và số bit lỗi $k$ trong một khối 7 bit tuân theo phân phối nhị thức
\begin{equation}
  \Pr(k \text{ lỗi}) = \binom{7}{k} p^k (1-p)^{7-k}.
\end{equation}
Vì mã (7,4) sửa đúng khi và chỉ khi có không quá một lỗi, xác suất một khối
được giải mã đúng là
\begin{equation}\label{eq:pkhoi}
  \Pr(\text{khối đúng}) = (1-p)^7 + 7p(1-p)^6 .
\end{equation}

\canviet{Mô tả cách \code{channel.py} cài đặt mô hình: tham số $p$, \code{seed},
có chế độ gây đúng $k$ lỗi hay không (cần cho mục~\ref{sec:5.2}); không sửa
đè danh sách đầu vào.}

% ---------------------------------------------------------------------
\section{Cách chạy chương trình}
\label{sec:4.5}

\canviet{(1) Yêu cầu: phiên bản Python, các file đặt chung thư mục.
(2) Lệnh chạy. (3) Dữ liệu nhập: chuỗi ký tự, $p$, \code{seed}.
(4) Một ví dụ chạy và kết quả in ra (nên dùng chuỗi \code{"A"} để khớp
Mục~\ref{sec:3.7}). Cần mã file chạy hoặc ảnh chụp kết quả thật.}

% ---------------------------------------------------------------------
\section{Một số xử lý đặc biệt}
\label{sec:4.6}

\canviet{Chỉ nêu những xử lý chương trình thực sự có, ví dụ: chuỗi rỗng; ký
tự ngoài ASCII (UTF-8, 2--4 byte mỗi ký tự); byte không hợp lệ sau khi giải
mã sai được thay bằng U+FFFD; kiểm tra kiểu và độ dài đầu vào
(\code{TypeError}/\code{ValueError}); các trường hợp biên của kênh ($p = 0$,
$p = 1$, $p \notin [0,1]$). Lưu ý: hội chứng bằng 0 chỉ cho biết $r$ là một từ
mã, không chứng tỏ không có lỗi; bộ giải mã (7,4) không biết số lỗi thật.}
